\section{HC Dynamics, Real Actions and the Geometry Interface}
\label{v480:hc-dynamics-geometry}

This section carries the conditional coupling program through its
projected dynamics, action and bilinear geometry. The retained source
provides the ordered master equation
(Theorem~\ref{thm:master-field-equation}), the curvature action
(Theorem~\ref{thm:variational-form}) and a conditional Einstein
reconstruction (Proposition~\ref{prop:bilinear-einstein-reconstruction}).
These are not already a common solved field system. In particular,
the source has not derived an identification between the microscopic
coordinate gamma matrices and the inverse of its reconstructed coframe.
The distinction becomes essential below.

All finite examples here use commuting effective amplitudes and declared
adjoints. A dagger denotes the ordinary positive internal adjoint where
specified, not an identification with the Lorentzian metric or the
native QAM comparison. Dimensionless local charts are used in the
explicit examples; a physical normalization requires the scale record
of Section~\ref{v480:scale-record}.

\subsection{Root Positivity Is an Additional Dynamical Obligation}

For a complex two-mode input $\chi=(u,v)$ set
\[
 X_0=|u|^2+|v|^2,\quad X_1=|u|^2-|v|^2,\quad
 X_2=2\operatorname{Re}(\overline u v),\quad
 X_3=2\operatorname{Im}(\overline u v).
\]
Then $X_0^2=X_1^2+X_2^2+X_3^2$. Let $U$ be the three-cycle
permutation matrix, and introduce the orthonormal Hermitian matrices
\[
 E_0=I/\sqrt3,\qquad E_1=(U+U^\dagger)/\sqrt6,\qquad
 E_2=i(U-U^\dagger)/\sqrt6.
\]
The following is an additional quadratic root readout, not an evaluated
native mass operator:
\[
 L_\zeta=X_0E_0+\zeta X_1E_1+\zeta X_2E_2,\qquad \zeta\geq0.
\]

\begin{proposition}[Balance and full-input positivity are different]
\label{v480:hc-root-domain}
For $X_0>0$ and $L_\zeta\geq0$, its candidate root ratio is
\[
 \frac{\operatorname{Tr}(L_\zeta^2)}{(\operatorname{Tr}L_\zeta)^2}
 =\frac{1+\zeta^2[1-(X_3/X_0)^2]}3.
\]
Positivity for every input $\chi$ holds if and only if
$\zeta\leq1/\sqrt2$. The aligned real or purely imaginary slice
has $X_3=0$; attaining $2/3$ on it requires $\zeta=1$, outside
the full-input positive range.
\end{proposition}
\begin{proof}
Orthogonality gives $\operatorname{Tr}L_\zeta=\sqrt3X_0$ and
$\operatorname{Tr}L_\zeta^2=X_0^2+\zeta^2(X_1^2+X_2^2)$.
In the three Fourier modes, the traceless part has eigenvalues
$\zeta\sqrt{2/3}(X_1\cos\varphi_j-X_2\sin\varphi_j)$,
where $\varphi_j=2\pi j/3$. Their minimum over all inputs with
fixed $X_0$ is $-\zeta\sqrt{2/3}X_0$, attained on the equator
$X_3=0$. Comparing with $X_0/\sqrt3$ proves the bound.
\end{proof}

Thus an isometric coefficient choice does not derive a physical
positive domain or a mass law. On the aligned slice
$\chi=ib(\cos(\theta/2),\sin(\theta/2))$, $b>0$, the
$\zeta=1$ positive domain consists of the three closed chambers
$|\theta-2\pi j/3|\leq\pi/12$, modulo $2\pi$.
Preparation and preservation of a chosen chamber must be checked.
A fixed-radius phase minimum of a separately constrained energy does
not prove stationarity or stability against amplitude, hidden-mode
and background variations. Taking absolute values of signed roots
does not preserve the displayed ratio in general.

\subsection{Projection Retains Hidden States and Moving Frames}

Write $D_\mu=\partial_\mu+\mathcal B_\mu$ and suppose
$\Gamma^0$ is invertible. The ordered master equation gives the
formal generator
\[
 \partial_t\Psi=\mathcal G\Psi,\qquad
 \mathcal G=-\mathcal B_0-(\Gamma^0)^{-1}
       \sum_{a=1}^3\Gamma^a(\partial_a+\mathcal B_a)
       -i(\Gamma^0)^{-1}\Omega.
\]
This rearrangement asserts neither a self-adjoint Hamiltonian nor
well-posedness on an unspecified domain. Its full HC contribution is
\[
 \mathcal G_{\rm HC}=-(\Gamma^0)^{-1}
       [\kappa\Gamma^\mu n_\mu+i\gamma I]H_C.
\]
The order matters. Positivity of $H_C$ alone is not a decay theorem.

\begin{proposition}[Exact projected autonomy criterion]
\label{v480:hc-projected-autonomy}
Let $u$ be the realification of a finite modal state,
$\dot u=\mathcal G_{\mathbb R}u$, and let $r=T(t)u$ be a
surjective real-linear extraction. For any right inverse $R_s$,
write $u=R_sr+z$, $Tz=0$. Then
\[
 \dot r=K r+q_z,\quad
 K=(\dot T+T\mathcal G_{\mathbb R})R_s,\quad
 q_z=(\dot T+T\mathcal G_{\mathbb R})z.
\]
The instantaneous equation is independent of every hidden lift if and
only if $(\dot T+T\mathcal G_{\mathbb R})\ker T=0$.
\end{proposition}
\begin{proof}
Differentiate $r=Tu$ and substitute the decomposition. Two lifts of
the same $r$ differ by an arbitrary element of $\ker T$, giving
both necessity and sufficiency.
\end{proof}

For example, $\dot f=-iHf$ with $H=S+iB$ Hermitian has
$S$ real symmetric and $B$ real skew. Writing $f=x+ir$ gives
$\dot x=Bx+Sr$, $\dot r=Br-Sx$. For constant blocks, eliminating
$x$ yields the exact memory equation
\[
 \dot r(t)=Br(t)-Se^{Bt}x(0)
  -\int_0^t Se^{B(t-s)}S r(s)\,ds.
\]
It is not automatically a Markovian damping law. In two real
dimensions an autonomous linear generator gives
$\dot\theta=2(\omega+c\cos\theta+d\sin\theta)$ on the aligned
slice. Forward invariance requires inward signs at the lower and upper
boundaries. Summing over the equally spaced lower and upper boundary
triples forces $\omega=0$; the same signs then force each boundary
value to vanish. A nonzero first harmonic cannot have those six
distinct zeros. Only scalar radial evolution remains. This does not
exclude nonlinear feedback or one separately selected chamber.

\subsection{A Conservative Inertia Can Be Derived Without Viscosity}

For a declared positive canonical two-mode action, let
\begin{align*}
 L&=\frac i2(f^\dagger\dot f-\dot f^\dagger f)-\mathcal E(f),\\
 e&=(\cos(\theta/2),\sin(\theta/2)),\qquad
 N=\begin{pmatrix}0&-1\\1&0\end{pmatrix},\\
 f&=(a+ib)e+cNe,\qquad b>0.
\end{align*}
The four real coordinates are essential. The kinetic term is
$b\dot a-a\dot b-bc\dot\theta$, or
$-2a\dot b-bc\dot\theta$ up to a total derivative.
Setting $a=c=0$ before variation deletes the canonical partners.

\begin{proposition}[Regular hidden-block elimination]
\label{v480:hc-collective-inertia}
Put $q=(b,\theta)$, $h=(a,c)$ and $D=\operatorname{diag}(2,b)$.
If $\mathcal E=V+j^Th+\tfrac12h^TWh$, with $W$ positive
definite, algebraic elimination in the action gives
\[
 h=-W^{-1}(D\dot q+j),\qquad
 L_{\rm eff}=\tfrac12(D\dot q+j)^TW^{-1}(D\dot q+j)-V.
\]
The induced kinetic metric is $DW^{-1}D$; its coordinate dependence
must be retained in the Euler--Lagrange equations.
\end{proposition}
\begin{proof}
Modulo the total derivative, $L=-h^TD\dot q-\mathcal E$.
Its $h$ equation is $Wh+D\dot q+j=0$. Substitution completes
the square and proves the formula.
\end{proof}

For $H=mI+\omega\sigma_y$, $m>|\omega|$, and
$\mathcal E=f^\dagger Hf$, this gives
\[
 L_{\rm eff}=\dot b^2/m+
      b^2(\dot\theta-2\omega)^2/(4m)-mb^2,
 \qquad I_\theta=b^2/(2m).
\]
The equation $\ddot\theta=-2(\dot b/b)(\dot\theta-2\omega)$
describes reversible amplitude exchange: the apparent damping changes
sign under contraction. It is not viscosity. The native positive
modal adjoint and coefficients have not been extracted. More generally,
a full-action pullback tests tangent equations only; omitted normal
field equations remain necessary.

\subsection{The Null Kinetic Term and the Real-Action Choice}

Suppress a constant commuting closure background and take flat fixed
spacetime, constant real null $n$, $A=0$, and $H_C=hT_0$, with
$T_0$ fixed Hermitian, $c_T=\operatorname{Tr}T_0^2>0$.
For $\mathcal B_\mu=\kappa H_Cn_\mu$ the printed curvature
action reduces exactly to
\[
 F_{\mu\nu}=\kappa T_0(n_\nu\partial_\mu h-n_\mu\partial_\nu h),
 \qquad
 L_{\rm HC}=\frac{\kappa^2c_T}{2}(n\mathbin{\cdot}\partial h)^2.
\]
Before imposing nullity its principal tensor is
$\kappa^2c_T(n^\mu n^\nu-n^2g^{\mu\nu})$; for nonzero null
$n$ and $\kappa\ne0$ it has rank one. A positive homogeneous
temporal inertia is not a Lorentzian wave operator or fluid viscosity.
If minimal exact closure makes $H_C=0$ identically on a branch, this
contribution vanishes there. Vanishing at one point is insufficient.

There is a separate reality issue. For ordinary Hermitian internal
$A_\mu,H$, set
\[
 D_\mu^AH=\partial_\mu H+i[A_\mu,H],\quad
 F^A_{\mu\nu}=\partial_\mu A_\nu-\partial_\nu A_\mu+i[A_\mu,A_\nu],
 \quad X_{\mu\nu}=n_\nu D_\mu^AH-n_\mu D_\nu^AH.
\]
Then $F=iF^A+\kappa X$. The literal $-\operatorname{Tr}(FF)/4$
can be complex, with real part
\[
 \operatorname{Re}L_F=\tfrac14\operatorname{Tr}(F^AF^A)
                         -\tfrac{\kappa^2}4\operatorname{Tr}(XX).
\]
The electric kinetic signs are opposite for signature $(-,+,+,+)$.
This is a convention-dependent diagnostic, not a native ghost theorem.
An explicitly \emph{different} unitary-internal comparison action is
\[
 L_F^\dagger=-\tfrac14\operatorname{Tr}(F_{\mu\nu}^\dagger F^{\mu\nu})
 =-\tfrac14\operatorname{Tr}(F^AF^A)
  -\tfrac{\kappa^2}4\operatorname{Tr}(XX).
\]
It is not silently substituted into Theorem~\ref{thm:variational-form}.
The traced cross commutator vanishes, but its field equations must be
derived from its own allowed variations.

\subsection{Probe Reality, Matrix Currents and Background Reactions}

Assume a factorization into internal and spin spaces,
$\bar\Psi=\Psi^\dagger\eta_D$, $\eta_D^\dagger=\eta_D$, constant Hermitian
$K^\mu=\eta_D\Gamma^\mu$, and
$(\eta_D\Omega)^\dagger=\eta_D\Omega$. Internal Hermitian fields
commute with the spin weights. Then
\[
 \operatorname{Re}(i\kappa\bar\Psi\Gamma^\mu n_\mu H\Psi)=0.
\]
A genuinely real symmetric probe action therefore removes this ordered
transport term, while retaining a real $\gamma H$ mass-like coupling.
This is a change of probe law, not an equivalent reading of the original.
For noncommuting weights the Hermitian part is instead
$i[K^\mu n_\mu,H]/2$, not the original ordered product.

For the paired curvature action, $\Omega=mI+\gamma H$, and an
independent Hermitian matching target $\Lambda$, the real test action
with penalty $\lambda_H\operatorname{Tr}(H-\Lambda)^2$ has sources
\[
 j^\mu=\operatorname{Tr}_{\rm spin}
       [(I\otimes K^\mu)|\Psi\rangle\langle\Psi|],\qquad
 S=\operatorname{Tr}_{\rm spin}
       [(I\otimes\eta_D)|\Psi\rangle\langle\Psi|].
\]
Here $S$ need not be positive. Independent $A,H$ variations give
\begin{align*}
 E_A^\nu&=D_\mu^AF_A^{\mu\nu}
          +i\kappa^2n_\mu[H,X^{\mu\nu}]-j^\nu,\\
 E_H&=\kappa^2n_\nu D_\mu^AX^{\mu\nu}
           -2\lambda_H(H-\Lambda)-\gamma S.
\end{align*}
Integration by parts in the two curvature terms gives the displayed
derivatives; varying the probe and penalty terms gives the last terms.
The probe equation yields $D_\mu^Aj^\mu=-i\gamma[H,S]$ on shell,
so its ordinary trace is conserved. Gauge covariance also requires
the target reaction $E_\Lambda=2\lambda_H(H-\Lambda)$ in the
Ward identity. For fixed transforming $\Lambda$ and no other
reactions, all other equations on shell require $[H,\Lambda]=0$.
This is not an instruction to impose $E_\Lambda=0$ or $H=\Lambda$.
A functional $\Lambda[M,P]$ instead transfers the reaction through
its chain rule. The inherited scalar gradient-trace currents are not
automatically these matrix currents.

Keeping the original ordered equation gives a different balance:
for the same constant weights and suppressed closure load,
\[
 \partial_\mu(\Psi^\dagger K^\mu\Psi)
       =-2\kappa\Psi^\dagger(K^\mu n_\mu)H_C\Psi.
\]
Its sign is not fixed by $H_C\geq0$. An open completion needs an
explicit return law or reservoir; a moving comparison needs its own
compatibility equations. Neither is derived by calling this loss
gravitational transport. No native viscosity or reservoir is supplied.

\subsection{An Explicit Null Geometry, Kept Off Shell}

In the fixed Hermitian factorization above, with the closure load
suppressed, the inherited real coframe has the exact identity
\[
 E^a{}_\mu=\operatorname{Re}(\bar\Psi\Gamma^aD_\mu\Psi)
       =\tfrac12\partial_\mu j^a+\kappa v_H^a n_\mu,
 \quad j^a=\Psi^\dagger K^a\Psi,\quad
 v_H^a=\Psi^\dagger K^aH_C\Psi.
\]
Differentiation proves the first term; the purely internal Hermitian
gauge term has zero real part. Unlike the $i$-weighted probe-action
term, the real HC coframe contribution survives.

\begin{proposition}[Gradient flatness and a null-preserving update]
\label{v480:hc-null-coframe}
If $H_C=0$ and $E$ has rank four in this fixed-frame branch, then
$Y^a=j^a/2$ are local coordinates and $g=\eta_{ab}dY^a dY^b$
is flat. More generally let $E=E_0+\kappa v n^T$, with $E_0$
invertible and $d=1+\kappa n^TE_0^{-1}v\ne0$. For
$g=E^T\eta E$, $g_0=E_0^T\eta E_0$,
\[
 \det E=(\det E_0)d,\qquad
 n^Tg^{-1}n=\frac{n^Tg_0^{-1}n}{d^2}.
\]
Thus this nonsingular rank-one update preserves a null covector.
\end{proposition}
\begin{proof}
The inverse function theorem proves the flat-coordinate statement.
The matrix determinant lemma gives the determinant, and the rank-one
inverse formula gives $n^TE^{-1}=n^TE_0^{-1}/d$. Contracting with
$\eta^{-1}$ proves the second identity.
\end{proof}

Variable Clifford/adjoint weights or a noncommuting geometric spin
connection can escape the gradient statement; it is not a theorem
about every native or coarse-grained metric. HC alone has rank at
most one and cannot furnish a four-dimensional coframe.

For a concrete local branch use $\eta=\operatorname{diag}(-1,1,1,1)$,
$n=(-1,0,0,1)^T$, $l=\eta^{-1}n$, and
\[
 E=I+\kappa h(x,y)ln^T,\qquad
 g=\eta+2\kappa h\,nn^T,\qquad \det E=1.
\]
This dictionary can be realized by a common local probe, not just by
independent choices of $E_0,v_H$. In a declared three-internal-mode,
four-spinor representation take $K^a=(I,\alpha_i)$,
$\eta_D=\beta=\operatorname{diag}(I_2,-I_2)$,
$\Gamma^a=\beta K^a$, with Pauli off-diagonal $\alpha_i$.
Let $\Psi_1=(1,0,1,0)^T/\sqrt2$, $\Psi_3=0$,
$H_C=\operatorname{diag}(h,0,0)$, and set
$J^a=2(x^a+c^a)-l^a$, $c=(3,0,0,0)$.
On a future-timelike patch for $J$, put
\[
 m_J=\sqrt{(J^0)^2-|\boldsymbol J|^2},\quad
 \Psi_2=\sqrt{\frac{J^0+m_J}{2}}
 \begin{pmatrix}\xi\\
 (\boldsymbol\sigma\mathbin{\cdot}\boldsymbol J)\xi/(J^0+m_J)
 \end{pmatrix},\qquad \xi=(1,0)^T.
\]
Direct Pauli contraction gives $j^a=2(x^a+c^a)$ and $v_H=hl$.
This uses $\{\Gamma^a,\Gamma^b\}=-2\eta^{ab}I$ and is not
an identification of the native septonic carrier. It is off shell.

The null coordinates $u=z-t$, $v=(z+t)/2$ put this metric in the
form $2du\,dv+dx^2+dy^2+2\kappa h\,du^2$. Its nonzero curvature
components are determined by
\[
 R_{\mu i\nu j}=-\kappa h_{,ij}n_\mu n_\nu,\quad
 R_{\mu\nu}=-\kappa\Delta_\perp h\,n_\mu n_\nu,\quad R=0,
 \qquad i,j\in\{x,y\}.
\]
These follow directly from the Levi--Civita connection; a vanishing
scalar curvature does not imply flatness.

Varying an independent unconstrained metric in the paired HC action
at fixed $n,F$ \emph{before} this null restriction gives
\[
 T^{\rm HC}_{\mu\nu}=\kappa^2c_T|\nabla_\perp h|^2n_\mu n_\nu.
\]
The contracted curvature action vanishes on the profile, but its metric
variation need not. Omitting all other stresses and the cosmological
term, taking $\kappa\ne0$ and $\kappa_{\rm eff}$ as an input, the isolated equation
$G_{\mu\nu}=\kappa_{\rm eff}T^{\rm HC}_{\mu\nu}$ reduces to
\[
 \Delta_\perp h+a|\nabla_\perp h|^2=0,\quad
 a=\kappa_{\rm eff}\kappa c_T,\qquad
 \Delta_\perp e^{ah}=0.
\]
For $a\ne0$, any positive harmonic $F$ gives the local solution
$h=\log F/a$. Positive affine $F$ supplies nonconstant examples;
restrict further to $h\geq0$ if $H_C\geq0$ is required.
This is an isolated off-shell geometry/stress match, not a full probe,
gauge, projector or minimal-closure solution, nor a derivation of
Newton's constant. A metric-enforced null constraint or an induced
metric variation changes the variational problem. Restricting the
action to this null family first loses normal Einstein equations.

\subsection{A Necessary Current Test and a Same-Frame Exclusion}

For the constructed transverse profile let $Y=x+c$. Its induced
coordinate current and volume satisfy
\[
 \mathcal J^\mu=(E^{-1})^\mu{}_a j^a
       =2Y^\mu-2\kappa h(n\mathbin{\cdot}Y)l^\mu,
 \qquad \sqrt{|\det g|}=1,\qquad
 \nabla_\mu\mathcal J^\mu=8.
\]
The last equality uses $n\mathbin{\cdot}l=0$ and
$l\mathbin{\cdot}\partial h=0$. Hence the unchanged construction
cannot solve the closed real probe theory with this induced phase
current. A purely internal Hermitian gauge field does not alter this
coframe or cancel a central divergence by a commutator. The number
$8$ is a normalization-dependent dimensionless diagnostic, not a
predicted production rate.

There is a stronger pointwise test which does not require the
factorized current construction.

\begin{proposition}[Conditional same-frame trace obstruction]
\label{v480:hc-same-frame-obstruction}
Suppose the same probe, adjoint and derivative define a real
classical full-rank coframe
\[
 E^a{}_\mu=s_0\operatorname{Re}(\bar\Psi\Gamma^aD_\mu\Psi),
 \qquad s_0\in\mathbb R\setminus\{0\},
\]
and make the additional identification
$\Gamma^\mu=(E^{-1})^\mu{}_a\Gamma^a$ in the same ordered
master equation. Set $R_\Psi=(i\Gamma^\mu D_\mu-\Omega)\Psi$
and $w=\bar\Psi\Omega\Psi$. Then
\[
 \operatorname{Im}(\bar\Psi R_\Psi)=4/s_0-\operatorname{Im}w.
\]
In particular, a real $w$ excludes an on-shell full-rank coframe
under these same-frame assumptions.
\end{proposition}
\begin{proof}
Reality of the inverse frame and its contraction with $E$ give
\[
 4=(E^{-1})^\mu{}_aE^a{}_\mu
  =s_0\operatorname{Re}(\bar\Psi\Gamma^\mu D_\mu\Psi).
\]
Multiply the master residual by the same adjoint and take the imaginary
part, using $\operatorname{Im}(iz)=\operatorname{Re}z$.
On shell the left side vanishes. This proves both statements.
\end{proof}

A sufficient reality condition is
$(\eta_D\Omega)^\dagger=\eta_D\Omega$; Hermiticity of $\Omega$
alone is insufficient without its relation to the adjoint. For finite
commuting amplitudes the same identity gives
\[
 \|R_\Psi\|\geq
 \frac{|4/s_0-\operatorname{Im}w|}{\|\eta_D\Psi\|}.
\]
Neither a finite real normalization nor a change of connection removes
the obstruction while all its assumptions are retained. It uses no
weak-field or Clifford approximation.

This is \emph{not} a universal inconsistency theorem for the paper.
The retained reconstruction has not derived that additional same-frame
identification. A reference/microscopic probe and a separately derived
effective geometry are not excluded. Changing the coframe functional,
deriving a changed probe equation from a composite-metric action, or
allowing a nonreal bilinear changes the problem and requires a new
common action and current balance. Merely inserting an imaginary
constant to meet the trace condition is not a derivation.

\subsection{Native Identification Still Required}
\label{v480:hc-geometry-obligations}

The construction establishes conditional inertia, a null kinetic
tensor, explicit real-action alternatives and a local curved metric
with an isolated stress match. It also identifies a definite forbidden
direct identification. It does not establish masses, an irreversible
fluid, a photon rest frame, a wormhole or inter-sector information flow.
The retained variational and Einstein statements must be read with
these qualifications, not promoted by a successful compiler run.

A native continuation must fix the common carrier and adjoint, decide
independent versus functional HC variations, derive every current and
background reaction from that same action, and specify the relation
between microscopic gamma matrices and effective geometry. Only then
can a common solution be tested against all tangent and normal
equations, the positive root domain and the physical scale record.
The executable HC root-readout, projected-residual, collective-action,
null-axis, real-current and bilinear-geometry audits in the repository
retain expanded examples and counterexamples. Their finite regression
checks support these identities; they are not native closure or
proof-assistant certificates.
